\chapter{Un enfant des routes}
\label{chap:un-enfant-des-routes}

À des lieues de là, peut-être même à l'autre bout du pays, une caravane s'était
arrêtée en pleine journée.

Les chariots avaient quitté la route pour se ranger sur le bas-côté. Les
chevaux attendaient, dételés pour certains, tandis que les voyageurs
profitaient de cette halte imprévue. Pourtant, personne ne semblait vraiment
pressé de repartir.

Soudain, un cri déchira le calme.

Il venait de l'un des chariots, celui de la famille Shinzawa. Quelque chose
d'inhabituel s'y déroulait, un événement suffisamment important pour avoir
interrompu le voyage de toute la caravane.

La naissance de leur premier fils.

\scenebreak

Les Shinzawa étaient des marchands itinérants. Ils parcouraient les routes avec
la caravane, transportant dans leur chariot des étoffes de toutes sortes qu'ils
vendaient au gré des villes et des villages traversés. Tissus simples pour les
vêtements du quotidien, étoffes épaisses pour les saisons froides, pièces plus
colorées ou plus délicates : leur marchandise variait selon ce qu'ils avaient
pu acheter, échanger ou vendre au cours de leurs voyages.

Ils possédaient également quelques talents de couturiers. Lorsque l'occasion se
présentait, ils réparaient un vêtement, reprenaient une couture ou ajustaient
une pièce pour un client avant de reprendre la route avec les autres voyageurs.

Et ce fut justement dans leur chariot, au milieu des étoffes qui accompagnaient
leur vie sur les routes, que naquit Kamatsu.

Enveloppé dans les étoffes de ses parents, Kamatsu lança son premier cri au
monde.